%%%%%%%%%%  scan idx 292  |  folio 279  %%%%%%%%%%
Les
deux dernières catégories sont des sous-catégories pleines de la catégorie \(\Mod({}_A X)\) des \(A\)-Modules à gauche, les deux premières des sous-catégories triangulées pleines de \(D({}_A X)\) (\(= D(\Mod({}_A X))\)).

\textbf{Définition 2.2.} Soit \((X,A)\) un topos annelé. Sauf mention du contraire, nous poserons
\[
 \begin{gathered}
  K_\bullet({}_A X)=k(D^b({}_A X)_{\coh}),\\
  K^{\bullet}({}_A X)=k(D({}_A X)_{\uparf}),
 \end{gathered}
\]
où \(k(\ )\) est défini en (1.1), et
\[
 \begin{gathered}
  K_\bullet({}_A X)_{\naif}=k(\Coh({}_A X)),\\
  K^{\bullet}({}_A X)_{\naif}=k(\Loclib({}_A X)),
 \end{gathered}
\]
où \(k(\ )\) est défini par (1.5.2), la catégorie abélienne ambiante étant \(\Mod({}_A X)\)\footnote{On se permettra d'omettre l'indice \(A\) quand il n'en pourra résulter de confusion. Parfois aussi on laissera tomber l'indice «~naïf~», mais en le signalant explicitement.}. Nous dirons que \(K_\bullet({}_A X)\) (resp. \(K^{\bullet}({}_A X)\)) est le groupe de Grothendieck des complexes pseudo-cohérents (resp. parfaits) de \(A\)-Modules à gauche, et que \(K_\bullet({}_A X)_{\naif}\) (resp. \(K^{\bullet}({}_A X)_{\naif}\)) est le groupe de Grothendieck des \(A\)-Modules à gauche cohérents (resp. localement libres).

\textbf{Remarques 2.3.}
\begin{enumerate}[label=\alph*)]
\item On peut définir des variantes de \(K_\bullet(\ )\) et \(K^{\bullet}(\ )\) en posant
\[
 K_\bullet^{\loc}({}_A X)=k(D^{\loc b}({}_A X)_{\coh})
\]
(où «~\(\loc b\)~» signifie «~à cohomologie localement bornée~»), et
\[
 K^{\bullet}_{\loc}({}_A X)=k(D({}_A X)_{\parf}).
\]
Ces groupes coïncident avec ceux définis en (2.2) lorsque l'objet final de \(X\) est quasi-compact, mais d'une manière générale se prêtent moins bien au calcul des intersections.
